\documentclass[11pt]{article}
\usepackage{amsmath,amssymb,amsthm}
\newtheorem{theorem}{Theorem}
\newtheorem{lemma}{Lemma}
\title{Invariant Costs of One-Pass Butterfly Permutations}
\author{}\date{}\begin{document}\maketitle

\paragraph{Status.}
This note proves a finite combinatorial classification and an exact
reciprocal identity. Neither is a superconstant-saving compiler. In
particular, the final two-denominator realization has a linear
feature-space limitation stated explicitly below. These are written
proofs with exhaustive small exact checks, not Lean declarations.

\paragraph{Notation.}
Put $H=2^d$. At layer $k=0,\ldots,d-1$ of the butterfly, independently
swap or preserve each pair of wires differing only in bit $k$.
Let $\pi$ be the resulting permutation, with all $Hd/2$ switches
independent and uniform. A cost matrix $C=(C_{hy})$ is invariant if
$S_C(\pi)=\sum_h C_{h,\pi(h)}$ is independent of all switch settings.
The input Haar vectors are
\[
 u_{i,g}(h)=\mathbf1_{\{h\mathbin{\mathrm{div}}2^{i+1}=g\}}
                   (1-2h_i),\qquad
 0\leq i<d,\quad 0\leq g<H/2^{i+1}.
\]
The output Haar vectors use the opposite bit order:
\[
 v_{j,s}(y)=\mathbf1_{\{y\bmod2^j=s\}}(1-2y_j),\qquad
 0\leq j<d,\quad 0\leq s<2^j.
\]
Together with the constant vector, each family is an orthogonal basis
of $\mathbb R^H$.

\begin{lemma}[Two-path probabilities]
Fix distinct inputs $h,h'$ and distinct outputs $y,y'$. Let $k$ be
the highest bit where the inputs differ, and $l$ the lowest bit where
the outputs differ. Then
\[
 \Pr(\pi(h)=y,\ \pi(h')=y')=
 H^{-2}\begin{cases}1&l<k,\\2&l=k,\\0&l>k.\end{cases}       \tag{1}
\]
Also $\Pr(\pi(h)=y)=H^{-1}$.
\end{lemma}
\begin{proof}
A specified input-output pair has a unique butterfly path, which
prescribes $d$ switch outcomes. If $l<k$, the two paths use disjoint
switches and prescribe $2d$ outcomes. If $l=k$, they share exactly
one switch, with compatible opposite exits, and prescribe $2d-1$
outcomes. If $l>k$, their prescriptions at the layer-$k$ switch
require both inputs to use the same exit, which is impossible.
\end{proof}

\begin{theorem}[Complete invariant-cost space]
The invariant costs are exactly the span of row-only costs,
column-only costs, and the product Haar costs
$u_{i,g}(h)v_{j,s}(y)$ with $i>j$. This space has dimension
\[
                         H+\frac{Hd}{2}.                    \tag{2}
\]
More precisely, on the orthogonal product-Haar basis, the covariance
operator of the assignment matrix is diagonal. On the product with
input level $i$ and output level $j$, its eigenvalue is
\[
 \lambda_{ij}=\begin{cases}
 0&i>j,\\2/H&i=j,\\1/H&i<j.
 \end{cases}                                                 \tag{3}
\]
The eigenvalue is relative to the squared Frobenius norm of the cost
matrix. Row-only and column-only costs have variance zero.
\end{theorem}
\begin{proof}
Let $A_k(h,h')$ indicate that the highest differing input bit is $k$,
and let $B_l(y,y')$ indicate that the lowest differing output bit is
$l$. On the space of product Haar costs, the second-moment operator
given by (1) is
\[
 \frac1H I+\frac1{H^2}\sum_{k=0}^{d-1}
 A_k\otimes\left(\sum_{l<k}B_l+2B_k\right).             \tag{4}
\]
The mean of every product Haar statistic is zero. The vector
$u_{i,g}$ is an eigenvector of $A_k$, with eigenvalues $2^k$ for
$k<i$, $-2^i$ for $k=i$, and zero for $k>i$.
Similarly $v_{j,s}$ is an eigenvector of $B_l$, with eigenvalues
$2^{d-l-1}$ for $l>j$, $-2^{d-j-1}$ for $l=j$, and zero for $l<j$.
Substitution in (4) gives (3). In particular distinct product Haar
costs have zero covariance. A cost is invariant precisely when its
variance is zero, since every switch setting has positive probability.
This proves the claimed span. Its dimension is
\[
 2H-1+\sum_{i=1}^{d-1}\frac{H}{2^{i+1}}
                      \sum_{j=0}^{i-1}2^j
 =H+\frac{Hd}{2}.
\]
The same formula gives dimension one when $d=0$.
\end{proof}

\paragraph{Equivalent cut invariants.}
For every $k=0,\ldots,d$, every input block fixing the bits at least
$k$, and every output prefix fixing the bits below $k$, exactly one
routed wire belongs to both. The corresponding rectangle costs are
therefore invariant. Their span is the space in the theorem: their
product-Haar expansions give precisely the modes $i>j$, together
with row and column costs. The full-permutation invariant space only
contains row and column costs, so its classification cannot be applied
unchanged to the butterfly family.

\begin{lemma}[Exact two-denominator realization]
Let $D_0,D_1$ be distinct numbers larger than all distinct pool poles
$p_y$. Choose arbitrary $q_y$, a linear polynomial $L(t)=ut+v$,
and residues
\[
 a_y=q_y(D_0-p_y)(D_1-p_y)+L(p_y).
\]
For a binary input profile $\alpha_h$, set $D_h=D_0$ if
$\alpha_h=1$ and $D_h=D_1$ otherwise, and set $b_h=-L(D_h)$.
Then for every permutation $\pi$,
\[
 \sum_h\frac{a_{\pi(h)}+b_h}{D_h-p_{\pi(h)}}
 =\sum_y q_y(D_0-p_y)-Hu
       +(D_1-D_0)\sum_h\alpha_hq_{\pi(h)}.             \tag{5}
\]
Thus an invariant assignment cost of the displayed separated form
gives an exact common zero after subtraction of a constant bias.
\end{lemma}
\begin{proof}
For $D_h=D_0$ the summand equals $q_{\pi(h)}(D_1-p_{\pi(h)})-u$;
for $D_h=D_1$ it equals $q_{\pi(h)}(D_0-p_{\pi(h)})-u$.
Summing these identities gives (5).
\end{proof}

\paragraph{Why two denominator states alone do not give a saving.}
Suppose every query chooses, at each head, one of the same two fixed
head functions $f_{0,h}(\theta_h)$ and $f_{1,h}(\theta_h)$.
Writing its choice as $\alpha_{hj}\in\{0,1\}$, its total score is
\[
 \sum_h f_{1,h}(\theta_h)+
       \sum_h\alpha_{hj}
           (f_{0,h}(\theta_h)-f_{1,h}(\theta_h)).       \tag{6}
\]
Across queries this belongs to a fixed linear space of dimension at
most $H+1$. Fixed query biases produce an affine translate of that
space. Consequently this family cannot shatter more than $H+1$
queries. The limitation applies to (5), including arbitrary moving
row parameters and arbitrary permutation branches. It does not apply
to general query-dependent denominators, additional denominator
states, or higher-order constructions that genuinely leave (6).
Realizing useful parts of the larger invariant space with those
additional freedoms remains open.
\end{document}
